\section{Vincit omnia Veritas}
\label{sec:VincitOmniaVeritas}
\begin{quotationx}
\emph{Is the initiate preoccupied with exhibiting in forms that will astonish, amaze, or terrify ... }

On the contrary, if there is anything radically opposed to the style of the true initiate, it is precisely this behavior. By definition, the initiate is an occult being and his path is neither visible nor penetrable. He is elusive, not to be pigeonholed. He arrives from the direction contrary to that toward which all gazes are fixed and takes the most natural seeming vehicle for his supernatural action. He may be an intimate friend, companion, or lover; he may be sure of possessing all your heart and confidence. But he will always be something different, other than what he lets be known. We will perceive this `other' only when we've penetrated his domain. And then perhaps we will have the feeling of having been standing on an abyss.

Men desire that what they are be known, that what they do be acknowledged, and that we be pleased by the quality of their performance. In the words of Agrippa, we have learned how different the law is that governs the magus and the hermeticist. They judge all exhibitionism and personality to be puerile. There is no adept. He does not exist. He does not speak. They but seek to net the wind, who are diverted by such things. The hermeticist has reached a state that categorically avoids all reaction to human judgment. He has stopped taking an interest in what others may think of him, or say about him, just or unjust, good or bad. He knows only that certain things must happen: he provides the precise means and conditions for them and that is all. He does not pretend the action is his own. He is pure instrumentality. `Self-affirmation' is a mania he does not recognise. And the farther he advances, the more deeply his centre sinks into a superindividual and superpersonal range, like one of the great forces of Nature, while those on whom he acts have the impression of being free.
\flright{\textsc{Julius Evola}, \emph{The Hermetic Tradition}}
\end{quotationx}


\paragraph{The Masses}


In \emph{The Crisis of the Modern World}, \textbf{Rene Guenon} describes the qualities of the masses, the called, and the elect or elite. Regarding the \textbf{masses}, he writes:

\begin{quotex}
It is doubtless true that the masses have always been led in one way or another, and it could be said that their part in history consists primarily in allowing themselves to be led, since they represent a predominantly passive element ... but in order to lead them today it is sufficient to possess oneself of purely material means, taking the word matter this time in its ordinary sense, and this clearly shows to what depths the present age has sunk; and at the same time these same masses are made to believe that they are not being led, but that they are acting spontaneously and governing themselves, and the fact that they believe this to be true gives an idea of the extent of their unintelligence.
\end{quotex}


\paragraph{The Called}


The unintelligence of the \textbf{called} is more subtle. The Gospels say that ``many are called but few are chosen''. That is because the called have an interest yet possess only the possibilities of knowledge. Therefore, they are prone to various seductions that will lead them astray. As Guenon warns:

\begin{quotex}
Many efforts are thereby expended in vain, and thus it is that many of those who sincerely desire to react against the modern outlook have been reduced to impotence; having failed to discover the essential principles, without which all action is doomed to be ineffective, they have allowed themselves to be lured into blind alleys from which there is no possible way of escape.
\end{quotex}


This is quite unfortunate when sincere minds get closed due to an infatuation with those blind alleys. The modernistic spirit can be difficult to overcome. For example, egalitarianism, historicism, and psychologism are prevalent even among those who outward claim to be committed to Tradition. Guenon explains the sinister nature of this spirit:

\begin{quotex}
Without any doubt the modernistic spirit, which is truly diabolical in every sense of the word, is striving by all the means in its power to prevent these elements, today isolated and scattered, from achieving the cohesion that is necessary if they are to exert any real influence on the general mentality.
\end{quotex}


\paragraph{The Chosen}


Next is the \textbf{chosen}, or elect, who will constitute the elite. Again from the Gospels: ``For false Christs and false prophets shall arise, and shall show signs and wonders to seduce, if it were possible, even the elect.'' But the chosen are no longer able to be seduced in virtue of the inner realisation they have achieved. Because this inner realization is certain and secure, Guenon is far from pessimistic and despairing. He concludes:

\begin{quotex}
Those who will be successful in overcoming all these obstacles and in triumphing over the hostility of an environment opposed to all spirituality will doubtless be few in number; but, once again, it is not numbers that count here, for this is a realm where the laws are quite other than those of matter. There is therefore no occasion for despondency; and even were there no hope of achieving any visible result before the collapse of the modern world through a catastrophe, that would still not be a valid reason to refrain from embarking upon a work extending in scope far beyond the present time. Those who may feel tempted to give way to discouragement should remind themselves that nothing accomplished within the order can ever be lost, that confusion, error and darkness can enjoy no more than a specious and purely ephemeral triumph, that every kind of partial and transitory disequilibrium must perforce contribute towards the great equilibrium of the whole, and that nothing can ultimately prevail against the power of truth for their motto they should take the one adopted in former times by certain initiatory organisations in the West:
\end{quotex}


\textit{Vincit omnia Veritas}


\flrightit{Posted on 2010-04-14 by Cologero}

\begin{center}* * *\end{center}

\begin{footnotesize}
\begin{sffamily}
\texttt{GFJF on 2010-04-14 at 23:04 said:}

The passage from Evola called to mind one from Boethius which is shorter but in essence the same: ``For as often as a man receives the reward of fame for his boasting, the conscience that indulges in self congratulation loses something of its secret merit.'' Even more creeping, I have noticed, is silent self praise, when one behaves humbly but secretly thinks oneself superior, perhaps not even aware of the hypocrisy. Profane elitists succumb to this trap, one way or the other; and it's vital for us, when dealing in ideas like `the elect' and `the masses', to be careful we do not.

\hfill

\texttt{Mikkel on 2018-04-14 at 21:19 said:}

Aurelius as well noted that ``For that arrogance which grows in arrogance for its own lack of arrogance is the most irksome of all''. This is stated within a Evola book stating the differences between humility and modesty.

Work that is done with the heights in mind that correspond to the order of things will always be true and right, that is something that cannot be taken.

\hfill

\texttt{Alex on 2022-04-14 at 20:47 said:}

Very often, when reading Evola, I find myself wishing he would get to the point or simplify his message. This quote from his ``Hermetic Tradition'' really took my breath away. If I should ever again feel that he is confusing me I will look up to this quote that I pinned to my wall board and tell myself here is your master.

\end{sffamily}
\end{footnotesize}

